\section{Summary/Delta Synchronization}
\label{sec:sync}

Pairwise merge is sufficient for convergence (\cref{prop:convergence}) but on its own would require peers to ship full state to one another whenever they reconcile. For non-trivial contracts --- a long-running chat room, a reputation set with many writers --- this is prohibitive. The platform therefore extends the contract interface with three additional functions that let peers exchange only the difference between their states.

\subsection{Interface}

\begin{definition}[Sync interface]
\label{def:sync}
In addition to $(\statespace_\contract, \merge_\contract, e_\contract, \valid_\contract)$ of \cref{def:algebra}, a contract provides:
\begin{align*}
  \summarize_\contract &: \statespace_\contract \to \Sigma_\contract,\\
  \getdelta_\contract &: \statespace_\contract \times \Sigma_\contract \to \Delta_\contract,\\
  \applydelta_\contract &: \statespace_\contract \times \Delta_\contract \to \statespace_\contract,
\end{align*}
where $\Sigma_\contract$ and $\Delta_\contract$ are contract-defined byte-string types.
\end{definition}

The summary type $\Sigma_\contract$ is intended to be much smaller than the full state. For an append-only log it might be the set of record identifiers, or a Merkle root over them; for a counter it might be the counter's current value; for a large set it might be a Bloom filter.

\subsection{Required Property}

For the protocol below to be correct it suffices to require that $\getdelta_\contract$ produce something that, when applied, dominates both peers' starting states in the partial order induced by $\merge_\contract$:

\begin{property}[Sync soundness]
\label{prop:sync}
For all $\sigma_A, \sigma_B \in \statespace_\contract$, let $s = \summarize_\contract(\sigma_B)$ and $\delta = \getdelta_\contract(\sigma_A, s)$. Then
\[
  \applydelta_\contract(\sigma_B, \delta) \;\merge_\contract\; \sigma_A
  \;=\;
  \applydelta_\contract(\sigma_B, \delta).
\]
\end{property}

\noindent Equivalently, applying the delta brings $\sigma_B$ at least as far up the lattice as $\sigma_A$ would have if merged in full. The contract author is responsible for $\summarize$ and $\getdelta$ being informative enough to satisfy this property; the runtime does not verify it.

\subsection{Protocol}

Two peers $A$ and $B$ reconcile their states for a contract $\contract$ in a two-step exchange:

\begin{enumerate}[leftmargin=*, itemsep=2pt]
  \item $A$ computes $s_A = \summarize_\contract(\sigma_A)$ and sends $s_A$ to $B$.
  \item $B$ computes $\delta_B = \getdelta_\contract(\sigma_B, s_A)$ and sends $\delta_B$ to $A$. $A$ applies the delta: $\sigma_A \gets \applydelta_\contract(\sigma_A, \delta_B)$.
\end{enumerate}

The protocol is symmetric: $B$ runs the same exchange against $A$ to incorporate updates $A$ has that $B$ does not. In practice the two directions are interleaved over a single connection. Bandwidth is bounded by $|s_A| + |s_B| + |\delta_A| + |\delta_B|$, all of which are application-defined; for an append-only log the summary is a list of record identifiers and the delta is the (typically small) set of records the other peer is missing.

\subsection{Propagation in a Subscription Tree}
\label{sec:subscription-tree}

Once a contract is published, peers near $\ell(\contract)$ on the ring host its state. Other peers register as subscribers (\cref{sec:ops-summary}); a subscription forms a soft-state link from the subscriber back toward the contract's home location, and together these links form a tree rooted at $\ell(\contract)$. An update arriving at any node of the tree is merged into local state and then propagated outward along the tree, with each pairwise hop using the summary/delta protocol of the previous subsection.

Because $\merge_\contract$ is commutative, simultaneous updates injected at multiple points of the tree do not require coordination. They propagate along the edges, meet at internal nodes, and merge; the order in which different updates arrive at a given node does not affect its final state (\cref{prop:convergence}). This is the ``color mixing'' picture described informally in the Freenet manual: each update is a color, every node ends up the same blend, and the network never has to vote on which color wins.

\subsection{Relationship to Other Mechanisms}

\paragraph{Versus operation-based CRDTs.} An operation-based CRDT \citep{shapiro2011crdt} propagates each individual operation and relies on a causal-broadcast layer to ensure exactly-once, causally ordered delivery. Freenet's protocol propagates state, not operations, and tolerates arbitrary loss, reordering, and duplication --- the merge function does the work that causal broadcast would otherwise have to.

\paragraph{Versus anti-entropy with Merkle trees.} Distributed key-value stores commonly use Merkle-tree anti-entropy \citep{decandia2007dynamo} to discover differences between replicas. That mechanism is a special case of the protocol here: the summary is a hash tree and the delta is the subtree the other replica is missing. The Freenet interface lets each contract pick a summary tailored to its data structure, so an application that knows its writes are append-only can use a list of identifiers rather than a hash tree, paying a single round trip rather than $O(\log n)$ of them.

\paragraph{Versus rsync.} \texttt{rsync} \citep{tridgell1999rsync} reconciles files by exchanging block-level checksums; it is generic but oblivious to application structure. Summary/delta is the same shape of protocol lifted to a level where the application defines the granularity. A chat-room contract can let two replicas reconcile by exchanging the set of message identifiers, not the byte layout of the encoded chat log.
